\section{Even orders: all-odd supports}
\label{sec:allodd}

This section treats the supports at even order all of whose elements are odd,
the hardest supports for the $\sigma$-bump, because the clauses of
\autoref{def:bump} then split by parity (\autoref{sec:allodd-parity}).

\begin{theorem}[all-odd supports]
\label{thm:alloddall}
Let $k$ be even, and let $S\subseteq\{1,\dots,k-1\}$ be primitive and
palindromic, with every element odd and $|S|\ge3$. Suppose that
$(k,S)\ne(6,\{1,3,5\})$. Then $S$ carries a $\sigma$-bump.
\end{theorem}

The support $\{1,3,5\}$ at $k=6$ carries no $\sigma$-bump at all: condition
(B1) leaves only $\sigma\in\{1,3,5\}$ and the two free positions $2$ and $4$
of $Q$, and in each of the twelve cases (B2) fails. It is treated separately
in \autoref{prop:order6cert}.

Supports containing $1$ are handled in \autoref{sec:allodd-one} by small
explicit carriers. For the others, with $a=\min S\ge3$, the carrier is a
\emph{half-arc}: a free position $y$ is put in $Q$ when the phase
$\nu y$ modulo $2k$ (shifted by $t\in\{0,k\}$ if $y$ is odd) lies in the
lower half $[0,k)$, for a suitable odd multiplier $\nu$. At the index $i$ the
reads of free positions are then governed by the translate $\nu i+\nu S$ of
the phase set, and the clause at $i$ can fail only if that translate lies in
one half of the circle, apart from pins and deleted taps. At odd indices
there are no pins, and the phases of $S$, being symmetric under $z\mapsto
k-z$, are spread too evenly to fit into one half once $\nu$ separates two
elements of $S$ below $k/2$ (\autoref{sec:allodd-halfarc}). At even indices,
for a unit $\nu$, failures for both choices of $t$ would make $\nu S$ contain
a whole coset of a subgroup of $\Z/2k$, which is impossible
(\autoref{sec:allodd-even}); the one gap is the index $2a$, closed by a
three-point condition (T) on $\nu$, which a lemma on three-term
progressions, \autoref{thm:lemmaA}, supplies (\autoref{sec:allodd-e3}). The
supports $\{a,k/2,k-a\}$, with a single element below $k/2$, get an explicit
$\nu$ (\autoref{sec:allodd-three}), and seventeen supports at orders
$k\le42$ escape these arguments and are given a carrier in a table. On a
first reading one can skip the case-by-case proofs and the finite checks.

\subsection{The parity split}
\label{sec:allodd-parity}

Throughout this section $k$ is even, $S\subseteq\{1,\dots,k-1\}$ is
palindromic, and reads are described in the vocabulary of
\autoref{sec:nao-reading}, in particular \eqref{eq:read-reads} and
\eqref{eq:read-pins}.

Suppose now that every element of $S$ is odd. Then at an odd index every tap
reads an even position, and at an even index an odd one. Two consequences are
used repeatedly. First, the clauses at odd indices involve only the even
elements of $Q$, and those at even indices only the odd ones. Second, at an
odd index no position is a pin, and since $\sigma$ is odd the only deleted tap
is $k-i$ (when it lies in $S$), which reads the position $k$. At an even index
no position is $\equiv0\pmod k$, and the deleted taps are those reading
$\sigma$ or $\sigma+k$.

\subsection{Supports containing 1}
\label{sec:allodd-one}

When $1\in S$ the taps $1$ and $k-1$ read the two neighbours of every
position, and very small carriers suffice.

\begin{lemma}[a three-point carrier]
\label{lem:allodd-mid}
Let $k$ be even and $S$ palindromic and all-odd, with $1\in S$, and let $b\in S$
with $5\le b\le k/2$. Then $\sigma=b$ and $Q=\{b-1,b,b+1\}$ form a $\sigma$-bump
for $S$.
\end{lemma}

\begin{proof}
(B1) holds because $b\pm1$ are even and $0\notin Q$. Note that $k\ge2b\ge10$,
and that $1,b,k-b,k-1\in S$; only these four taps are used.

\emph{Even indices.} Here every position is odd, and the only odd element of
$Q$ is $b=\sigma$. So a surviving unwrapped tap reads $0$ and a surviving
wrapped tap reads $1$. For a $1$: the tap $k-1$ reads $i+k-1$, which is wrapped
and is deleted only if $i+k-1=k+b$, that is $i=b+1$; at $i=b+1$ the tap $k-b$
reads $k+1$, wrapped and not deleted. For a $0$: the tap $1$ reads $i+1<k$,
deleted only if $i+1=b$; at $i=b-1$ the tap $b$ reads $2b-1<k$, not deleted.

\emph{Odd indices.} Here every position is even and the only deleted position
is $k$. For a $1$: at $i\ge3$ the tap $k-1$ reads the wrapped position
$i+k-1$, of residue $i-1$, which reads $1$ unless $i-1=b\pm1$. The remaining
cases are $i=1$, where the tap $b$ reads $b+1\in Q$; $i=b$, where the tap $1$
reads $b+1$; and $i=b+2$, where the tap $k-b$ reads $k+2$, of residue
$2\notin Q$ (as $b-1\ge4$). For a $0$: unless $i=k-1$ the tap $1$ reads the
unwrapped position $i+1$, which reads $0$ unless $i+1=b\pm1$. The remaining
cases are $i=k-1$ and $i=b$, where the taps $b$ and $k-1$ respectively read
$k+b-1$, of residue $b-1\in Q$; and $i=b-2$, where the tap $b$ reads
$2b-2<k$, not in $Q$ because $b\ge5$.
\end{proof}

\begin{lemma}[a carrier when $3\in S$]
\label{lem:allodd-three4}
Let $k\ge8$ be even and $S$ palindromic and all-odd, with $1,3\in S$. Then
$\sigma=3$ and $Q=\{2,3,6\}$ form a $\sigma$-bump for $S$.
\end{lemma}

\begin{proof}
Since $k\ge8$, $Q\subseteq\{1,\dots,k-1\}$, and (B1) holds because $2$ and $6$
are even. Only the taps $1$, $3$, $k-3$ and $k-1$ are used.

\emph{Even indices.} As in \autoref{lem:allodd-mid}, every position is odd and
the only odd element of $Q$ is $\sigma$, so a surviving tap reads $0$ if
unwrapped and $1$ if wrapped. For a $1$: the tap $k-1$ reads the wrapped
position $i+k-1$, of residue $i-1$, which is deleted only at $i=4$; there the
tap $k-3$ reads $k+1$, wrapped and not deleted. For a $0$: the tap $1$ reads
$i+1<k$, deleted only at $i=2$; there the tap $3$ reads $5<k$, not deleted.

\emph{Odd indices.} Every position is even, the only deleted position is
$k$, and the even elements of $Q$ are $2$ and $6$. So a surviving tap reads
$1$ exactly when it is unwrapped with residue in $\{2,6\}$, or wrapped with
residue outside $\{2,6\}$. For a $1$: at $i=1$ the tap $1$ reads $2$; at
$i=3$ the tap $3$ reads $6$; at $i=7$ the tap $k-3$ reads $k+4$, of residue
$4$; and at every other odd index, which has $i\ge5$ and $i\ne7$, the tap
$k-1$ reads $k+i-1$, of residue $i-1\notin\{0,2,6\}$. For a $0$: at $i=1$ the
tap $3$ reads $4$; at $i=5$ and at $i=k-1$ the taps $k-3$ and $3$
respectively read $k+2$, of residue $2$; and at every other odd index, which
has $i\le k-3$ and $i\ne1,5$, the tap $1$ reads $i+1<k$, of residue
$i+1\notin\{0,2,6\}$.
\end{proof}

\begin{proposition}[all-odd supports containing $1$]
\label{prop:allodd-one}
Let $k$ be even and $S\subseteq\{1,\dots,k-1\}$ palindromic and all-odd, with
$1\in S$, $|S|\ge3$ and $(k,S)\ne(6,\{1,3,5\})$. Then $S$ carries a
$\sigma$-bump.
\end{proposition}

\begin{proof}
If $3\in S$, then $k\ge8$: for $k\le6$ the odd elements of $\{1,\dots,k-1\}$
are among $1,3,5$, so $|S|\ge3$ would force the excluded support
$(6,\{1,3,5\})$. So \autoref{lem:allodd-three4} applies. If $3\notin S$, then
$k-3\notin S$ either, by palindromy, and $S$ has an element other than $1$ and
$k-1$; that element or its mirror is an odd $b\in S$ with $5\le b\le k/2$, and
\autoref{lem:allodd-mid} applies.
\end{proof}

\subsection{Half-arc carriers and the odd indices}
\label{sec:allodd-halfarc}

From here on $1\notin S$. Write $a=\min S$, so that $a\ge3$ and $\sigma=a$
throughout, and write $L=\{s\in S:2s<k\}$ for the lower half of $S$. Since $S$
is palindromic, $S=L\cup(k-L)$, together with $k/2$ when $k/2\in S$; and
$a\in L$ unless $S=\{k/2\}$.

For integers $\nu$, $y$ and $z$ write
\begin{equation}
H(z)=\mathbf 1[\,z\bmod 2k<k\,],
\qquad
H_\nu(y)=H(\nu y),
\label{eq:allodd-eps}
\end{equation}
where $z\bmod2k$ denotes the representative in $\{0,\dots,2k-1\}$. So
$H_\nu(y)$ records in which half of $\Z/2k$ the \emph{phase} $\nu y$
lies. Two elementary facts about the phases are used throughout. If $\nu$ is
odd then $\nu k\equiv k\pmod{2k}$, so for every integer $z$ and $y$
\begin{equation}
H(z+k)=1-H(z),\qquad H_\nu(y+k)=1-H_\nu(y).
\label{eq:allodd-flip}
\end{equation}
If also $y$ is odd, then $\nu y$ is odd, so $\nu y\not\equiv0,k\pmod{2k}$, and
\begin{equation}
H_\nu(-y)=1-H_\nu(y),\qquad H_\nu(k-y)=H_\nu(y),
\label{eq:allodd-sym}
\end{equation}
the second because $\nu(k-y)\equiv k-\nu y$, and $z\mapsto k-z$ maps each of the
open halves $U=\{1,\dots,k-1\}$ and $V=\{k+1,\dots,2k-1\}$ of $\Z/2k$ to
itself.

\begin{definition}[half-arc carrier]
\label{def:allodd-halfarc}
Let $\nu$ be odd and $t\in\{0,k\}$, and for an integer $y$ put $t_y=0$ if $y$
is even and $t_y=t$ if $y$ is odd. The \emph{half-arc carrier} $Q(\nu;t)$ of $S$
consists of $a$ together with every $y\in\{1,\dots,k-1\}\setminus S$ with
$H(\nu y+t_y)=1$.
\end{definition}

So the even free positions follow the half-arc of $\nu$, and the odd ones the
same half-arc shifted by $t$. By construction $Q\cap(\{0\}\cup S)=\{a\}$, so
(B1) holds at $\sigma=a$. The point of the construction is that a tap reading
a free position reads the half-arc pattern, wrapped or not.

\begin{lemma}[free reads]
\label{lem:allodd-free}
Let $Q=Q(\nu;t)$ with $\nu$ odd, and let $0<p<2k$ be a free position. Then
$p\in P$ if and only if $H(\nu p+t_p)=1$.
\end{lemma}

\begin{proof}
If $p<k$ this is the definition, since $p\ne a$. If $p>k$, put $p'=p-k$, which
has the parity of $p$ and lies outside $S\cup\{0\}$. Then $p\in P$ iff
$p'\notin Q$ iff $H(\nu p'+t_p)=0$ iff $H(\nu p'+t_p+k)=1$, by \eqref{eq:allodd-flip},
and $\nu p'+k\equiv \nu p\pmod{2k}$.
\end{proof}

The clauses at odd indices read only even positions, and all of these are free
except the deleted position $k$. They are settled by the following lemma,
whose idea is that the phases of $S$ are spread round the circle $\Z/2k$ too
evenly to fit into one half, even after one of them is removed.

\begin{lemma}[odd indices]
\label{lem:allodd-oddhalf}
Let $S$ be all-odd, let $\nu$ be odd and $t\in\{0,k\}$, and suppose that
$H_\nu(u)\ne H_\nu(v)$ for some $u,v\in L$. Then, for
$\sigma=a$ and $Q=Q(\nu;t)$, the clause at every odd index holds.
\end{lemma}

\begin{proof}
Let $i$ be odd. As noted in \autoref{sec:allodd-parity}, the only deleted
tap is $k-i$ (if it lies in $S$), and every other tap $j$ reads a free even
position $p=i+j$, whose read is $H(\nu p)=H(\nu i+\nu j)$ by \autoref{lem:allodd-free}.
So the clause fails exactly when the phases $\nu j\bmod2k$ of the surviving taps
all lie in $\{-\nu i,\dots,-\nu i+k-1\}$ or all in $\{-\nu i+k,\dots,-\nu i+2k-1\}$.

First I locate the phases. They are odd, so each lies in $U$ or in $V$, and
by \eqref{eq:allodd-sym} the tap $k-j$ has phase $k-\nu j$, which lies in the
same half as $\nu j$. Since $H_\nu(u)\ne H_\nu(v)$ both halves
occur. Let $x_-$ be the least phase in $U$, attained at the tap $u_0$, and
$y_-$ the least in $V$, attained at $v_0$. The taps $k-u_0$ and $k-v_0$ have
the phases $k-x_-\in U$ and $3k-y_-\in V$, so $2x_-\le k$ and $2y_-\le3k$.
Going round the circle, each of the four phases $x_-\le k-x_-<y_-\le3k-y_-$
is followed by the next at distances $k-2x_-$, $x_-+y_--k$, $3k-2y_-$ and
$x_-+y_--k$ respectively. Each distance is at most $k$, and equals $k$ only
if $x_-=k/2$ and $y_-=3k/2$; so at most $k-1$ consecutive residues avoid all
four phases, and $k-1$ only in that case. If $m$ consecutive residues contain
all four phases, the other $2k-m$ residues are consecutive and avoid them, so
$2k-m\le k-1$. Hence:
\begin{itemize}
\item[(a)] no $k$ consecutive residues contain all four phases;
\item[(b)] if $k+1$ consecutive residues contain them, then $x_-=k/2$ and
$y_-=3k/2$, and then every phase is $k/2$ or $3k/2$ (a phase $z\in U$ has
$z\ge x_-$ and $k-z\ge x_-$, and similarly in $V$).
\end{itemize}

If $k-i\notin S$, nothing is deleted and the clause holds by (a). Suppose
instead that $s_0=k-i\in S$. Then $\nu i\equiv k-\nu s_0$, and by
\eqref{eq:allodd-flip} the tap $j\ne s_0$ reads $1-H(\nu j-\nu s_0)$. If these reads
were all $0$, every phase would lie in the $k$ consecutive residues
$\{\nu s_0,\dots,\nu s_0+k-1\}$, which (a) excludes. If they were all $1$,
every phase other than $\nu s_0$ would lie in
$\{\nu s_0+k,\dots,\nu s_0+2k-1\}$, and with $\nu s_0$ itself every phase would lie in
the $k+1$ consecutive residues $\{\nu s_0+k,\dots,\nu s_0+2k\}$; so (b) gives
every phase in $\{k/2,3k/2\}$, and then the tap $j\ne s_0$ reads $0$ when
$\nu j\equiv \nu s_0$ and $1$ otherwise. The taps $u$ and $k-u$ are distinct (as
$2u<k$) and share one phase; $v$ and $k-v$ share the other. At most one of the
four taps is $s_0$, so the pair with the phase of $s_0$ contains a surviving
tap reading $0$, and the other pair a surviving tap reading $1$. This
contradicts the assumption that the reads were all $1$.
\end{proof}

\subsection{Supports with one element below \texorpdfstring{$k/2$}{k/2}}
\label{sec:allodd-three}

If $1\notin S$ and $L=\{a\}$, then $|S|\ge3$ forces $S=\{a,k/2,k-a\}$, with
$k/2$ odd. Here \autoref{lem:allodd-oddhalf} does not apply, since $L$ has a
single element, and a specific multiplier is used instead.

\begin{theorem}[three-element supports]
\label{thm:allodd-threeall}
Let $k=2h$ with $h$ odd, and $S=\{a,h,k-a\}$ with $3\le a<h$ odd and
$\gcd(a,h)=1$. Let $\beta\in\{2,4\}$ with $\beta\equiv a+h\pmod4$, suppose
that $3\beta<h$ (that is, $h\ge7$ if $\beta=2$ and $h\ge13$ if $\beta=4$), and
let $\nu\equiv a^{-1}(3h+\beta)\pmod{2k}$. Then $\sigma=a$ and $Q=Q(\nu;0)$ form a
$\sigma$-bump for $S$.
\end{theorem}

\begin{proof}
All residues in this proof are modulo $2k=4h$. Since $a$ is odd and prime to
$h$, it is a unit, so $\nu$ is well defined and $\nu a\equiv3h+\beta$. As
$a^2\equiv1\pmod 4$ and $3h+\beta\equiv3h+a+h\equiv a\pmod4$, one has
$\nu\equiv a(3h+\beta)\equiv1\pmod4$; so $\nu$ is odd and $\nu h\equiv h$. The phases
of the three taps are therefore
\[
\nu a\equiv3h+\beta,\qquad \nu h\equiv h,\qquad \nu(k-a)\equiv k-\nu a\equiv3h-\beta .
\]
(B1) holds by construction. By \autoref{lem:allodd-free}, a tap reading a free
position $p$ reads $H(\nu p)=H(\nu i+\nu j)$, where $j$ is the tap.

\emph{Deletions and pins.} The deleted positions are $k$, $a$ and $k+a$; the
pins are $h$ and $k-a$, which read $0$, and $k+h$ and $2k-a$, which read $1$
(the pins $a$ and $k+a$ are deleted). Solving $i+j=p$ for $j\in S$, $p$ in
this list and $0<i<k$ shows that a tap is deleted or reads a pin only at the
seven indices
\[
a,\quad h,\quad k-a,\quad h-a,\quad k-2a,\quad h+a,\quad 2a.
\]

\emph{Generic indices.} At any other index all three taps read free positions,
with phases $\nu i+\{3h+\beta,\,h,\,3h-\beta\}$. Going round the circle, the
three points $h<3h-\beta<3h+\beta$ cut it into arcs of lengths $2h-\beta$,
$2\beta$ and $2h-\beta$, all less than $2h$, so no half
$\{c,\dots,c+2h-1\}$ contains all three phases, and the reads are not all
equal.

\emph{The special indices.} At each of the seven special indices one names
a tap reading $0$ and a tap reading $1$, using the rule that the phase of a
position is the sum of the phases of the index and the tap, and that the
position reads $1$ exactly when its phase lies in $\{0,\dots,2h-1\}$. The
seven cases, and the check that every tap named survives, are in
\autoref{app:checks-threeall}.
\end{proof}

\subsection{The even indices with a unit multiplier}
\label{sec:allodd-even}

Now let $\nu$ be a unit modulo $2k$. Multiplication by $\nu$ is then a bijection
of $\Z/2k$, and it is convenient to work entirely with phases. Put
\[
\Phi=\{\nu s\bmod 2k: s\in S\},\qquad \alpha=\nu a\bmod2k .
\]
The elements of $\Phi$ are odd, and $\Phi=k-\Phi$ by \eqref{eq:allodd-sym}. Moreover
$H_\nu$ is non-constant on $S$ exactly when $\Phi$ meets both $U$ and $V$.

Let $t\in\{0,k\}$ and $Q=Q(\nu;t)$. Let $0<i<k$ be even, and put
$\phi=\nu i\bmod2k$.
Then $\phi$ is even, and $\phi\notin\{0,k\}$ because $\nu$ is a unit and
$i\not\equiv0,k\pmod{2k}$; the same holds for $\phi+k$. The tap $j$ reads the
odd position $p=i+j$, of phase $\nu p\equiv\phi+\nu j$. Since $p>j\ge a$, the tap is
deleted exactly when $p=a+k$, that is, when the phase of $p$ is $\alpha+k$.
Because $p\mapsto \nu p$ is a bijection, $p$ is an unwrapped pin iff $\nu p\in \Phi$,
and a wrapped pin iff $\nu p\in \Phi+k$. So a surviving tap whose position has
phase $z$ reads
\begin{equation}
0 \text{ if } z\in \Phi;\qquad 1 \text{ if } z\in \Phi+k;\qquad
\text{otherwise } H(z+t)=\begin{cases}\mathbf 1[z\in U],&t=0,\\
\mathbf 1[z\in V],&t=k,\end{cases}
\label{eq:allodd-evenread}
\end{equation}
by \autoref{lem:allodd-free}; here $z$ is odd, so it lies in $U$ or in $V$.
The phases read at $i$ form the \emph{read set} $\phi+\Phi$, less the point
$\alpha+k$ when that point is the phase of a deleted tap. For a residue $\tau$
consider the two conditions
\[
C_U(\tau):\ \ (\tau+\Phi)\cap U\subseteq \Phi,\qquad
C_V(\tau):\ \ (\tau+\Phi)\cap V\subseteq \Phi.
\]

\begin{lemma}[failing read sets]
\label{lem:allodd-windows}
With the notation above, suppose that the clause at the even index $i$ fails.
\begin{itemize}
\item[(i)] If every surviving tap reads $1$, then $C_U(\phi+k)$ holds when
$t=0$, and $C_V(\phi+k)$ when $t=k$.
\item[(ii)] If every surviving tap reads $0$, then $C_U(\phi)$ holds when
$t=0$, and $C_V(\phi)$ when $t=k$, unless some tap is deleted at $i$ and
$\alpha+k$ lies in $U$ (for $t=0$), respectively in $V$ (for $t=k$).
\item[(iii)] If some tap is deleted at $i$ and $i\ne2a$, then some surviving
tap reads $1$.
\end{itemize}
\end{lemma}

\begin{proof}
(i) Take $t=0$. Let $x\in \Phi$ with $z=\phi+x\in V$. If $z$ is the deleted point
then $z=\alpha+k\in \Phi+k$. Otherwise $z$ reads $1$; by
\eqref{eq:allodd-evenread} it is not in $\Phi$, and since $z\notin U$ it must lie
in $\Phi+k$. So $(\phi+\Phi)\cap V\subseteq \Phi+k$, and adding $k$ gives
$C_U(\phi+k)$. The case $t=k$ is the same with $U$ and $V$ exchanged.

(ii) Take $t=0$. Let $x\in \Phi$ with $z=\phi+x\in U$, and suppose $z$ is not
the deleted point; this is automatic unless a tap is deleted and
$\alpha+k\in U$. Then $z$ reads $0$, so $z\notin \Phi+k$, and as $z\in U$
\eqref{eq:allodd-evenread} forces $z\in \Phi$. The case $t=k$ is the same.

(iii) The deleted tap is $j_0=a+k-i$. The tap $k-a\in S$ reads
$i+k-a=k+(k-j_0)$, a wrapped pin, since $k-j_0\in S$. It is deleted only if
$k-j_0=a$, that is $i=2a$. So for $i\ne2a$ it reads $1$.
\end{proof}

Two failures of opposite type, one in $U$ and one in $V$, are incompatible. The
reason is that together they force $\Phi$ to contain a whole coset of a
subgroup, which a set of the form $\nu S$ cannot do.

\begin{lemma}[no coset]
\label{lem:allodd-nocoset}
Let $\nu$ be a unit modulo $2k$ and $S\subseteq\{1,\dots,k-1\}$. Then $\Phi=\nu S$
contains no coset $c+g\Z/2k\Z$ with $g\mid2k$ and $0<g<2k$.
\end{lemma}

\begin{proof}
Multiplication by $\nu^{-1}$ is an automorphism of $\Z/2k$, so it maps such a
coset onto a coset $s_0+g\Z/2k\Z$, which would then lie in $S$ (read modulo
$2k$). Since $g\mid 2k$ and $g<2k$, $g\le k$. The integers $s_0,s_0+g,
s_0+2g,\dots$ start below $k$ and increase in steps of at most $k$, so one of
them lies in $\{k,\dots,2k-1\}$; its residue is in the coset but not in $S$.
\end{proof}

\begin{lemma}[two failures force a coset]
\label{lem:allodd-cosetfail}
Let $k$ be even, and let $\Phi$ be a set of odd residues modulo $2k$ with
$\Phi=k-\Phi$ that meets both $U$ and $V$. Let $\phi,\phi'$ be even residues not in
$\{0,k\}$, and suppose that $C_U(\phi)$ and $C_V(\phi')$ hold. Then $\Phi$
contains a coset $c+g\Z/2k\Z$ with $g\mid 2k$ and $0<g<k$.
\end{lemma}

\begin{proof}
\emph{Reductions.} $C_U(\phi)$ implies $C_U(-\phi)$: if $x\in \Phi$ and
$z=x-\phi\in U$, then $k-x\in \Phi$ and $\phi+(k-x)=k-z\in U$, so $k-z\in \Phi$ and
$z\in \Phi$. Likewise $C_V(\phi')$ implies $C_V(-\phi')$, since $z\mapsto k-z$
also maps $V$ to itself. So I may take $0<\phi,\phi'<k$. The set $\Phi+k$ again
consists of odd residues, satisfies $\Phi+k=k-(\Phi+k)$ and meets $U$ and $V$, and
$C_U(\tau)$ for $\Phi$ is $C_V(\tau)$ for $\Phi+k$ and conversely. So passing
from $\Phi$ to $\Phi+k$ swaps $C_U$ and $C_V$: the hypotheses $C_U(\phi)$ and
$C_V(\phi')$ for $\Phi$ become $C_U(\phi')$ and $C_V(\phi)$ for $\Phi+k$, which
are the hypotheses of the lemma for $\Phi+k$ with $\phi$ and $\phi'$
interchanged. A coset in $\Phi+k$ gives one in $\Phi$, by adding $k$. So if
$\phi>\phi'$ I may replace $\Phi$ by $\Phi+k$, and I may therefore also take
$\phi\le\phi'$.

\emph{Translation into $(0,k)$.} Put $Y_+=\Phi\cap U$ and $Y_-=\{z:z+k\in \Phi\cap
V\}$, two nonempty sets of odd integers in $(0,k)$. By $\Phi=k-\Phi$ both are
closed under $z\mapsto k-z$ (for $Y_-$: $k-(m+k)\equiv(k-m)+k$). The
hypotheses give four rules:
\begin{itemize}
\item[(G1)] if $p\in Y_+$ and $p+\phi<k$, then $p+\phi\in Y_+$ (by $C_U(\phi)$);
\item[(G2)] if $m\in Y_-$ and $m<\phi$, then $\phi-m\in Y_+$ (the point
$2k-m\equiv k-(m+k)$ of $\Phi$ moves to $\phi-m\in U$);
\item[(G3)] if $m\in Y_-$ and $m+\phi'<k$, then $m+\phi'\in Y_-$ (by $C_V(\phi')$);
\item[(G4)] if $p\in Y_+$ and $p<\phi'$, then $\phi'-p\in Y_-$ (the point $k-p$ of
$\Phi$ moves to $k+\phi'-p\in V$).
\end{itemize}
Call a set $\Pi$ of integers in $(0,k)$ \emph{$n$-periodic} if $p\in\Pi$,
$0<z<k$ and $z\equiv p\pmod n$ imply $z\in\Pi$. The proof uses one
mechanism twice.

\emph{The engine.} Let $\Pi$ be $n$-periodic, $n>0$, and suppose that for
every $p\in\Pi$ there is $p'\in\Pi$ with $p'\equiv p+f\pmod n$. Then $\Pi$ is
$\gcd(n,f)$-periodic. Indeed, write $g=\gcd(n,f)$ and choose $y\ge0$ with
$yf\equiv g\pmod n$. If $p\in\Pi$, $0<z<k$ and $z\equiv p\pmod g$, then
$z\equiv p+jg\pmod n$ for some $j\ge0$; applying the hypothesis $jy$ times
gives $p''\in\Pi$ with $p''\equiv p+jyf\equiv z\pmod n$, and so $z\in\Pi$.

\emph{Step 1: $Y_+$ is $\gamma$-periodic, where $\gamma=\gcd(\phi,\phi')$.} By
(G1) and the symmetry of $Y_+$ ($p-\phi=k-((k-p)+\phi)$), $Y_+$ is closed under
$\pm\phi$ inside $(0,k)$, so it is $\phi$-periodic. Let $p\in Y_+$ and
$p_0=p\bmod\phi$; it is odd, so $0<p_0<\phi$ and $p_0\in Y_+$. Let $r\ge0$ be
largest with $p_0+r\phi<\phi'$. By (G1), $p_0+r\phi\in Y_+$, and by (G4),
$m=\phi'-p_0-r\phi\in Y_-$. Now $0<m\le\phi$, and $m$ is odd, so $m<\phi$, and
(G2) gives $\phi-m\in Y_+$. Since $(\phi-m)+\phi'=p_0+(r+1)\phi\equiv p
\pmod\phi$, every $p\in Y_+$ has some $p'\in Y_+$ with $p'+\phi'\equiv p\pmod\phi$.
Applying this to $k-p$ and reflecting, every $p\in Y_+$ has $p''=k-p'\in Y_+$ with
$p''\equiv p+\phi'\pmod\phi$. The engine with $n=\phi$, $f=\phi'$ proves Step
1.

\emph{Step 2: if $0<z<k$ is odd and $z+p\equiv0\pmod\gamma$ for some $p\in Y_+$,
then $z\in Y_-$.} Let $z_0=z\bmod\phi'$, which is odd, so $0<z_0<\phi'$. Then
$\phi'-z_0\equiv-z\equiv p\pmod\gamma$, so $\phi'-z_0\in Y_+$ by Step 1, and
(G4) gives $z_0\in Y_-$. Applying (G3) repeatedly, $z=z_0+q\phi'\in Y_-$.

\emph{Step 3: $Y_+$ is $g$-periodic, where $g=\gcd(\gamma,2k)$.} Let $p\in Y_+$,
and let $z=(k+p)\bmod\gamma$; it is odd, so $0<z<\gamma\le\phi$. Since
$(k-z)+p\equiv0\pmod\gamma$, Step 2 gives $k-z\in Y_-$, hence $z\in Y_-$, and
(G2) gives $\phi-z\in Y_+$, hence $k-\phi+z\in Y_+$. As $\gamma\mid\phi$,
$k-\phi+z\equiv k+z\equiv p+2k\pmod\gamma$. The engine with $n=\gamma$,
$f=2k$ proves Step 3.

\emph{Step 4: the coset.} The number $g$ is even, divides $2k$, and
$0<g\le\phi<k$. Fix $c\in Y_+$ and let $w$ be any residue $\equiv c\pmod g$,
represented in $[0,2k)$; it is odd. If $w<k$ then $w\in Y_+\subseteq \Phi$ by Step
3. If $w>k$, put $z=w-k$, so that $-z\equiv k-c\pmod g$. Let $z'=(-z)\bmod
\gamma$, which is odd with $0<z'<\gamma<k$. Since $z'\equiv k-c\pmod g$ and
$k-c\in Y_+$, Step 3 gives $z'\in Y_+$, and $z+z'\equiv0\pmod\gamma$, so $z\in Y_-$
by Step 2; that is, $w\in \Phi$. So $c+g\Z/2k\Z\subseteq \Phi$.
\end{proof}

The coset lemma disposes of every pair of failures except one: the all-$0$
failure of a read set that has lost its deleted point. By
\autoref{lem:allodd-windows}(iii) this can only happen at the index $2a$.
The following condition on $\nu$, which asks for an element of $L$ and a free
position read at the index $2a$, both in the half opposite to $a$, prevents it
there.

\begin{definition}[the three-point condition]
\label{def:allodd-e3}
Let $1\notin S$ and $a=\min S$. A unit $\nu$ modulo $2k$ \emph{satisfies the three-point
condition} (T) if
\begin{itemize}
\item[(Ta)] some $x\in L\setminus\{a\}$ has $H_\nu(x)\ne H_\nu(a)$,
and
\item[(Tb)] some $j\in S\setminus\{k-a\}$ with $(2a+j)\bmod k\notin S$ has
$H_\nu(2a+j)\ne H_\nu(a)$.
\end{itemize}
\end{definition}

\begin{proposition}[the half-arc carrier under (T)]
\label{prop:allodd-e3bump}
Let $S$ be all-odd with $1\notin S$, and let $\nu$ be a unit satisfying (T).
Then for $t=0$ or for $t=k$, $\sigma=a$ and $Q(\nu;t)$ form a $\sigma$-bump
for $S$.
\end{proposition}

\begin{proof}
(B1) holds by construction, and the clauses at odd indices hold for both
values of $t$ by \autoref{lem:allodd-oddhalf}, applied with $u=x$ from (Ta)
and $v=a$. It remains to show that for one of the two values of $t$ the
clause holds at every even index.

By (Ta), $2a<2x<k$, so $2a$ is an index, and $\Phi$ meets $U$ and $V$. Suppose
that both choices of $t$ fail, at the even indices $i_0$ (for $t=0$) and $i_1$
(for $t=k$).

I claim first that the failure at $i_0$ gives $C_U(\tau)$ for some even
$\tau\notin\{0,k\}$. By \autoref{lem:allodd-windows}(i) and (ii) this holds
with $\tau\in\{\nu i_0,\nu i_0+k\}$, unless every surviving tap at $i_0$ reads $0$,
some tap is deleted there, and $\alpha+k\in U$. In that case $i_0=2a$ by
\autoref{lem:allodd-windows}(iii). Moreover $\alpha\in V$, that is
$H_\nu(a)=0$. Take $j$ from (Tb). The tap $j$ is not deleted at $2a$,
since $2a+j\ne a+k$, and it reads a free position. By \autoref{lem:allodd-free}
its read is $H(\nu(2a+j))=H_\nu(2a+j)=1$, a contradiction. In the same
way the failure at $i_1$ gives $C_V(\tau')$ for some even
$\tau'\notin\{0,k\}$. Here the exceptional case has $\alpha+k\in V$, so
$H_\nu(a)=1$, and the tap $j$ reads $H(\nu(2a+j)+k)=
1-H_\nu(2a+j)=1$.

\autoref{lem:allodd-cosetfail} now puts a coset $c+g\Z/2k\Z$, $0<g<k$, inside
$\Phi$, and \autoref{lem:allodd-nocoset} says that there is none.
\end{proof}

\subsection{The three-point condition}
\label{sec:allodd-e3}

It remains to find a unit satisfying (T). For $k>90$ there are
two routes. Either some $x\in L\setminus\{a\}$ gives (T) through a unit that
witnesses the progression $2a-x$, $a$, $x$, provided $x<2a$ or
$x-2a\notin S$ (\autoref{lem:allodd-aproute}(i));
or $L$ is the chain $a,3a,\dots,(2m+1)a$, and \autoref{thm:lemmaA}, applied
to the rescaled triple $(k;1,2m+3)$, supplies the unit. Triples with
$\gcd(a,x,k)>1$ may have no witness, but these are rescaled copies of the $59$
failures at orders up to $90$ listed in \autoref{lem:allodd-pf} (see
\autoref{lem:allodd-lift}), and \autoref{lem:allodd-nofail} shows that they
cannot block both routes, except for $a=15$ at $k=120$, which is settled by
hand. The orders $k\le90$ are a finite check.

\begin{definition}[witnesses and admissible triples]
\label{def:allodd-adm}
Let $k$ be even and let $a,x$ be integers. A unit $\nu$ modulo $2k$
\emph{witnesses} the triple $(k;a,x)$ if
\begin{equation}
H_\nu(2a-x)=H_\nu(a)\ne H_\nu(x),
\label{eq:allodd-ap}
\end{equation}
and $(k;a,x)$ \emph{has a witness} if some unit does. The triple is
\emph{admissible} if $a<x$ are odd positive integers with $2x<k$, $x\ne3a$
and $\gcd(a,x,k)=1$. (The triple with $x=3a$ never has a witness, since then
$2a-x=-a$ and $H_\nu(-a)=1-H_\nu(a)$ for every odd $\nu$, by
\eqref{eq:allodd-sym}.)
\end{definition}

Note that $2a-x$, $a$ and $x$ form a three-term progression, so a witness
separates the progression after its middle term but not before it. The main
input of this section is \autoref{thm:lemmaA}, proved in
\autoref{sec:lemmaA}: \emph{every admissible triple $(k;a,x)$ with $k>90$ has
a witness}.

\begin{lemma}[the progression route]
\label{lem:allodd-aproute}
Let $S$ be primitive and all-odd with $1\notin S$.
\begin{itemize}
\item[(i)] Let $x\in L\setminus\{a\}$, and suppose that $x<2a$ or
$x-2a\notin S$. Then every unit $\nu$ witnessing $(k;a,x)$ satisfies (T).
\item[(ii)] Suppose that $|L|\ge2$ and that every $x\in L\setminus\{a\}$ has
$x>2a$ and $x-2a\in S$. Then $\gcd(a,k)=1$, and there is $m\ge1$ with
$2(2m+1)a<k$, $(2i+1)a\in S$ for $0\le i\le m$, and $(2m+3)a\notin S$.
\end{itemize}
\end{lemma}

\begin{proof}
(i) (Ta) holds with this $x$, since $H_\nu(x)\ne H_\nu(a)$.
For (Tb) take $j=k-x\in S$, which differs from $k-a$. If $x<2a$, then
$(2a+j)\bmod k=2a-x\in(0,a)$, which is not in $S$ by the minimality of $a$. If
$x>2a$ (equality is excluded by parity), then $2a+j=k-(x-2a)<k$, which lies in
$S$ only if $x-2a$ does. Finally, by \eqref{eq:allodd-flip},
$H_\nu(2a+j)=H_\nu(2a-x+k)=1-H_\nu(2a-x)=
1-H_\nu(a)$.

(ii) Each $x\in L\setminus\{a\}$ has $x-2a\in L$, and descending from $x$
by steps of $2a$ stays in $L$ until it reaches $a$. So every element of $L$
is $\equiv a\pmod{2a}$, in particular divisible by $a$. Let $f=\gcd(a,k)$,
which is odd. It divides every element of $L$, every $k-x$ with $x\in L$, and
$k/2$ (as $f$ is odd and divides $k$). So $f$ divides $\gcd(S\cup\{k\})=1$.
Let $(2m+1)a$ be the largest element of $L$; $m\ge1$ since $|L|\ge2$, and the
descent shows $(2i+1)a\in S$ for $0\le i\le m$. Suppose $y=(2m+3)a\in S$. Then
$2y<k$ contradicts the maximality, $2y=k$ gives $a\mid k$, and $2y>k$ puts
$k-y$ in $L$, so $a\mid k-y$ and again $a\mid k$. But $a\ge3$ and
$\gcd(a,k)=1$.
\end{proof}

Up to order $90$ some admissible triples have no witness, and these failures
propagate upwards by scaling.

\begin{lemma}[small orders]
\label{lem:allodd-pf}
Let $(n;a,x)$ be admissible with $n\le90$. Then $(n;a,x)$ has a witness
unless it is one of the $59$ triples
\begin{center}
\small
\begin{tabular}{r|l}
$n$ & $(a,x)$\\\hline
12 & (1,5) (3,5)\\
16 & (1,5) (1,7) (3,7)\\
20 & (1,7) (3,5) (7,9)\\
22 & (1,5) (1,9) (5,7)\\
24 & (1,5) (1,11) (3,7) (5,7) (5,9) (7,9)\\
28 & (1,7) (3,13) (7,9)\\
30 & (1,11) (1,13) (7,13)\\
36 & (1,7) (1,15) (1,17) (5,11) (5,13) (7,17) (11,15)\\
40 & (1,5) (1,15) (3,15) (5,17) (7,19) (9,13) (15,19)\\
42 & (1,5) (1,9) (1,17) (3,19) (5,7) (11,15) (11,19) (13,15)\\
48 & (3,11) (15,23)\\
60 & (1,23) (11,13) (13,27) (19,21)\\
70 & (1,17) (3,25) (9,13) (13,15)\\
90 & (1,33) (17,21) (19,41) (23,39)
\end{tabular}
\end{center}
Denote this list by $\mathcal F$.
\end{lemma}

\begin{proof}
This is a finite check: for every admissible triple with $n\le90$ not in
$\mathcal F$, a search over the units $\nu$ modulo $2n$ finds a witness.
\end{proof}

\begin{lemma}[lifting]
\label{lem:allodd-lift}
If $(n;a,x)$ has a witness, then so does $(gn;ga,gx)$ for every $g\ge1$.
Consequently, let $k>90$ and let $(k;a,x)$ satisfy every condition of
admissibility except perhaps $\gcd(a,x,k)=1$. If $(k;a,x)$ has no witness,
then, with $g=\gcd(a,x,k)$, the triple $(k/g;a/g,x/g)$ lies in $\mathcal F$.
\end{lemma}

\begin{proof}
Let $\nu$ be a unit modulo $2n$ witnessing $(n;a,x)$. Reduction
$\Z/2gn\to\Z/2n$ maps units onto units, so there is a unit $\nu'$ modulo $2gn$
with $\nu'\equiv \nu\pmod{2n}$. For every integer $y$,
$\nu'\,gy\bmod 2gn=g\,(\nu y\bmod2n)$, which is less than $gn$ exactly when
$\nu y\bmod 2n<n$. So $H_{\nu'}(gy)$ at order $gn$ equals $H_\nu(y)$
at order $n$, and $\nu'$ witnesses $(gn;ga,gx)$.

For the second part, $g$ divides the odd number $a$, so it is odd; hence
$k/g$ is even, and $(k/g;a/g,x/g)$ is admissible. It has no witness, by the
first part. If $k/g>90$ this contradicts \autoref{thm:lemmaA}, so $k/g\le90$,
and \autoref{lem:allodd-pf} puts the triple in $\mathcal F$.
\end{proof}

Above order $90$ the lifted failures can still occur, but only in a way that
the three-point condition can absorb.

\begin{lemma}[failures above order $90$]
\label{lem:allodd-nofail}
Let $k>90$ be even, and let $S$ be primitive and all-odd with $1\notin S$ and
$(k,a)\ne(120,15)$. Suppose that every $x\in L\setminus\{a\}$ for which $(k;a,x)$ has
a witness satisfies $x>2a$ and $x-2a\in S$; that is, route (i) of
\autoref{lem:allodd-aproute} is available for no $x$. Then $(k;a,x)$ has a witness for
every $x\in L\setminus\{a,3a\}$.
\end{lemma}

\begin{proof}
Suppose that $(k;a,x_0)$ has no witness for some $x_0\in L\setminus\{a,3a\}$;
I derive a contradiction. Let $\mathcal M$ be the set of orders $n$ such that
$(n;a',x')\in\mathcal F$ for some $a',x'$ with $n\mid k$ and $a=(k/n)a'$, and
let $\ell=\operatorname{lcm}\mathcal M$. For $y\in L\setminus\{a,3a\}$ write
$g_y=\gcd(a,y,k)$. If $(k;a,y)$ has no witness, then
\autoref{lem:allodd-lift} gives $(k/g_y;a/g_y,y/g_y)\in\mathcal F$, so
$k/g_y\in\mathcal M$. In particular $\mathcal M$ is nonempty (take
$y=x_0$). Every element of $\mathcal M$ divides $k$, so $\ell\mid k$; write
$k=\ell f$. For every $n\in\mathcal M$, $n\mid\ell$, so $f$ divides $k/n$;
hence $f\mid g_y$ whenever $(k;a,y)$ has no witness. In particular
$f\mid g_{x_0}\mid a$, so $f$ is odd.

I claim that $f$ divides every $y\in L$, by induction on $y$. For $y=a$ this
was just shown. For $y\ne a$ there are three cases. If $(k;a,y)$ has a
witness, then $y-2a\in L$ by hypothesis, $f\mid y-2a$ by induction, and so
$f\mid y$. If it
has none and $y=3a$, then $f\mid a\mid y$. If it has none and $y\ne3a$, then
$f\mid g_y\mid y$. Since $f$ is odd and divides $k$, it also divides $k/2$
and every $k-y$; so $f$ divides $\gcd(S\cup\{k\})=1$. Hence $\ell=k$.

It remains to see that $\ell=k$ is impossible. Each $n\in\mathcal M$ comes
from a failure $(n;a',x')\in\mathcal F$ with $a'/n=a/k$, so all the orders in
$\mathcal M$ carry failures with one and the same ratio $a'/n$. Inspecting
$\mathcal F$, a ratio occurs at two different orders only three times:
$3/12=7/28$, $3/24=5/40$ and $1/16=3/48$. So either $\mathcal M$ is a single
order, and $\ell\le90$, or $\ell$ divides $\operatorname{lcm}(12,28)=84$,
$\operatorname{lcm}(24,40)=120$ or $\operatorname{lcm}(16,48)=48$. As
$\ell=k>90$, this leaves $k=120$ with $a/k=1/8$, that is, $a=15$, which is
excluded.
\end{proof}

\begin{theorem}[the three-point condition]
\label{thm:allodd-e3all}
Let $S$ be primitive and all-odd, with $1\notin S$ and $|L|\ge2$, and not one
of the supports listed in \autoref{tab:allodd-base}. Then some unit $\nu$
modulo $2k$ satisfies (T).
\end{theorem}

\begin{proof}
\emph{Orders $k>90$.} If some $x\in L\setminus\{a\}$ has a witness
for $(k;a,x)$ and satisfies $x<2a$ or $x-2a\notin S$,
\autoref{lem:allodd-aproute}(i) gives (T). Otherwise, if
$(k,a)\ne(120,15)$, \autoref{lem:allodd-nofail} applies, and every $x\in L\setminus\{a,3a\}$ has a
witness, so it satisfies $x>2a$ and $x-2a\in S$; so does $x=3a$, since
$a\in S$. So \autoref{lem:allodd-aproute}(ii) applies. It gives
$\gcd(a,k)=1$ and $m\ge1$ with $(2i+1)a\in S$ for $i\le m$, $2(2m+1)a<k$
and $(2m+3)a\notin S$. As $a\ge3$, $2(2m+3)<6(2m+1)\le2(2m+1)a<k$, so the
triple $(k;1,2m+3)$ is admissible, and by \autoref{thm:lemmaA} some unit
$\nu_1$ witnesses it:
$H_{\nu_1}(-(2m+1))=H_{\nu_1}(1)\ne H_{\nu_1}(2m+3)$. By
\eqref{eq:allodd-sym}, $H_{\nu_1}(2m+1)=1-H_{\nu_1}(-(2m+1))\ne
H_{\nu_1}(1)$. Since $a$ is odd and prime to $k$, it has an inverse
$\bar a$ modulo $2k$. Put $\nu=\nu_1\bar a$, a unit, so that
$\nu\,ya\equiv\nu_1y\pmod{2k}$ and $H_\nu(ya)=H_{\nu_1}(y)$ for every
integer $y$. Let $j=(2m+1)a$, which lies in $L\setminus\{a\}$. Then (Ta)
holds with $x=j$, since
$H_\nu(j)=H_{\nu_1}(2m+1)\ne H_{\nu_1}(1)=H_\nu(a)$.
For (Tb), $j\ne k-a$ and $2a+j=(2m+3)a\le2(2m+1)a<k$, so
$(2a+j)\bmod k=(2m+3)a\notin S$, and
$H_\nu(2a+j)=H_{\nu_1}(2m+3)\ne H_{\nu_1}(1)=H_\nu(a)$.

\emph{The order $k=120$ with $a=15$.} Suppose again that route (i) of
\autoref{lem:allodd-aproute} is available for no $x$, so that every
$x\in L\setminus\{a\}$ with a witness has $x>30$ and $x-30\in S$. Let
$x\in L\setminus\{15,45\}$ have no witness. By \autoref{lem:allodd-lift},
$(120/g;15/g,x/g)\in\mathcal F$ with $g=\gcd(15,x,120)$, and $g\ne1$ since the
orders in $\mathcal F$ are at most $90$. The only entries of $\mathcal F$ of
the forms $(40;5,\cdot)$, $(24;3,\cdot)$ and $(8;1,\cdot)$ are $(40;5,17)$ and
$(24;3,7)$, so $x\in\{35,51\}$. Every element of $L$ therefore lies in
$\{15,35,45,51\}$, by induction: an $x\in L$ with a witness has $x-30\in L$,
and $x-30\in\{15,35,45,51\}$ with $x<60$ forces $x=45$. If $35\notin L$, then
$3$ divides every element of $L$, hence every element of $S$ (as $k/2=60$ is
even, $S=L\cup(k-L)$), contradicting primitivity; so $35\in L$, and in the same
way, with $5$ in place of $3$, $51\in L$. Now the unit $\nu=11$ satisfies (T),
with residues modulo $240$. First $H_{11}(15)=0$, as $11\cdot15=165$. (Ta)
holds with $x=51$, since $11\cdot51\equiv81$. For (Tb), if $45\notin S$ take
$j=15$, so that $2a+j=45\notin S$ and $11\cdot45\equiv15$; if $45\in S$ take
$j=k-35=85$, so that $2a+j=115\notin S$ (as $5<a$) and $11\cdot115\equiv65$.

\emph{Orders $k\le90$.} These are a finite check; since
$1\notin S$ and $|L|\ge2$, $k\ge12$. For each such $k$ it has three steps.
\begin{enumerate}
\item Let $F_k$ be the set of the pairs $(a,3a)$ with $6a<k$, which have no
witness by \eqref{eq:allodd-sym}, together with the pairs $(ga',gx')$ with
$g$ odd and $(k/g;a',x')\in\mathcal F$; for instance
$F_{44}=\{(1,3),(3,9),(5,15),(7,21)\}$. For every odd pair $a<x<k/2$ not in
$F_k$, a witness of $(k;a,x)$ is found by a search over the units $\nu$ modulo $2k$.
\item By \autoref{lem:allodd-aproute}(i), $S$ then satisfies (T) unless
every $x\in L\setminus\{a\}$ has $(a,x)\in F_k$ or $x-2a\in L$. The lower
halves $L$ with this property are enumerated, for each odd $a\ge3$ (as
$1\notin S$) with $2a<k$, starting from $L=\{a\}$ and adding the
candidates $x$ in increasing order. Each is completed to a support with and without $k/2$
(when $k/2$ is odd), and the supports that are not primitive or have
$|L|<2$ are discarded. This leaves $411$ supports.\footnote{Counts such as
this one, and the numbers of triples and orders quoted in
\autoref{sec:lemmaA} and \autoref{app:lemmaA-finite}, describe the size of
the finite checks. They are not themselves formalised; what is formalised is
that every case met by each check passes it (\autoref{app:lean}).}
\item For each of these $411$ supports except the twelve with $|L|\ge2$ in
\autoref{tab:allodd-base}, a unit $\nu$ satisfying (T) is recorded.
\end{enumerate}
\end{proof}

\subsection{Proof of \autoref{thm:alloddall}}
\label{sec:allodd-proof}

The seventeen supports that escape the arguments above are listed in
\autoref{tab:allodd-base}, each with a half-arc carrier. The multiplier $\nu$
is a residue modulo $2k$, and it need not be a unit. Indeed, for the supports
$\{3,9,13,19\}$ and $\{5,7,15,17\}$ at $k=22$ and $\{7,13,17,23\}$ at
$k=30$ there is no half-arc carrier with a unit multiplier: for every unit
$\nu$ modulo $2k$ and both $t\in\{0,k\}$, the pair $(\min S,Q(\nu;t))$ fails
to be a $\sigma$-bump. Several entries of the table use non-units; for these
three supports, with $\nu=11$, $11$ and $3$, no unit would do.

\begin{table}
\centering
\small
\begin{tabular}{rlrc@{\qquad}rlrc}
$k$ & $S$ & $\nu$ & $t$ & $k$ & $S$ & $\nu$ & $t$\\\hline
10 & $\{3,5,7\}$ & 3 & $0$ & 22 & $\{9,11,13\}$ & 7 & $0$\\
14 & $\{5,7,9\}$ & 5 & $k$ & 24 & $\{5,9,15,19\}$ & 3 & $k$\\
16 & $\{3,7,9,13\}$ & 3 & $0$ & 30 & $\{7,13,17,23\}$ & 3 & $0$\\
18 & $\{7,9,11\}$ & 5 & $k$ & 30 & $\{7,13,15,17,23\}$ & 3 & $0$\\
22 & $\{3,9,13,19\}$ & 11 & $k$ & 36 & $\{5,13,23,31\}$ & 3 & $k$\\
22 & $\{3,9,11,13,19\}$ & 3 & $0$ & 40 & $\{3,9,31,37\}$ & 5 & $k$\\
22 & $\{5,7,15,17\}$ & 11 & $0$ & 40 & $\{7,19,21,33\}$ & 3 & $0$\\
22 & $\{5,11,17\}$ & 5 & $k$ & 40 & $\{9,13,27,31\}$ & 11 & $0$\\
 & & & & 42 & $\{5,15,27,37\}$ & 3 & $k$
\end{tabular}
\caption{The supports of \autoref{lem:allodd-base}, with the parameters
$(\nu,t)$ of a half-arc carrier $Q(\nu;t)$ that, with $\sigma=\min S$, forms a
$\sigma$-bump. The multiplier $\nu$ is an odd residue modulo $2k$, not
necessarily a unit. Five have three elements, of the form $\{a,k/2,k-a\}$ with
$3\beta\ge k/2$ in the notation of \autoref{thm:allodd-threeall}, and the
other twelve have $|L|\ge2$.}
\label{tab:allodd-base}
\end{table}

\begin{lemma}[seventeen small supports]
\label{lem:allodd-base}
For each of the seventeen supports $S$ of \autoref{tab:allodd-base}, with the
listed $\nu$ and $t$, $\sigma=\min S$ and $Q(\nu;t)$ form a $\sigma$-bump.
\end{lemma}

\begin{proof}
This is a finite check of (B1) and (B2) for seventeen explicit pairs
$(\sigma,Q)$.
\end{proof}

\begin{proof}[Proof of \autoref{thm:alloddall}]
If $1\in S$, \autoref{prop:allodd-one} applies. Assume $1\notin S$, and let
$a=\min S\ge3$; then $a\in L$, since otherwise $S=\{k/2\}$. If $S$ is listed
in \autoref{tab:allodd-base}, \autoref{lem:allodd-base} applies. Otherwise
there are two cases.

\emph{$|L|\ge2$.} \autoref{thm:allodd-e3all} gives a unit $\nu$ satisfying
(T), and \autoref{prop:allodd-e3bump} gives the $\sigma$-bump $(a,Q(\nu;t))$
with $t\in\{0,k\}$.

\emph{$|L|=1$.} Then $L=\{a\}$, and since $|S|\ge3$, $S=\{a,k/2,k-a\}$ with
$h=k/2$ odd. Any common divisor of $a$ and $h$ divides $a$, $h$, $k-a$ and
$k$, so $\gcd(a,h)=1$ by primitivity, and $3\le a<h$. Let $\beta$ be as in
\autoref{thm:allodd-threeall}. If $3\beta\ge h$, then $h\le11$, and $(h,a)$ is
one of $(5,3)$, $(7,5)$, $(9,7)$, $(11,5)$ and $(11,9)$, all with $\beta=4$
(for $\beta=2$ the only candidate order, $h=5$, has no suitable $a$); so $S$
would be listed in \autoref{tab:allodd-base}. Hence $3\beta<h$, and
\autoref{thm:allodd-threeall} gives the $\sigma$-bump.
\end{proof}
